\chapter{Profilo utente}
\label{cap:profilo_utente}

In un approccio di progettazione centrata sull'utente (\textit{User-Centered Design}), è fondamentale definire non solo i requisiti tecnici, ma anche i bisogni, le abilità e le limitazioni di chi utilizzerà il sistema. Di seguito verranno illustrati i profili dei tre utenti individuati per la piattaforma \textbf{Guided Tours}: il Configuratore, il Volontario e il Fruitore.

\section{Configuratore (Amministratore del sistema)}
Il configuratore è un esponente dell'organizzazione che sovrintende alla pianificazione delle visite. È l'utente con il massimo potere sul prodotto e gestisce l'intero ciclo di vita dei tour, dalla creazione delle tipologie alla gestione dei luoghi.

\subsection{Caratteristiche psicologiche}
Il configuratore possiede uno stile cognitivo di tipo analitico e sistematico, necessario per coordinare risorse umane e logistiche. L'attitudine è estremamente positiva, essendo il successo dell'organizzazione legato all'efficienza dello strumento. La motivazione è alta, poiché il sistema è il mezzo principale per lo svolgimento della propria attività lavorativa.

\begin{table}[H]
\centering
\begin{tabular}{|l|l|}
\hline
\textbf{Caratteristica} & \textbf{Tipologia} \\ \hline
Stile cognitivo & Analitico, Verbale, Sistematico \\ \hline
Attitudine & Molto Positiva \\ \hline
Motivazione & Alta \\ \hline
\end{tabular}
\caption{Caratteristiche psicologiche del Configuratore.}
\label{tab:conf_psicologiche}
\end{table}

\subsection{Caratteristiche di conoscenza}
Trattandosi di un ruolo di responsabilità, si suppone un livello di istruzione medio-alto. L'alfabetizzazione informatica deve essere consolidata per gestire pannelli amministrativi complessi.

\begin{table}[H]
\centering
\begin{tabular}{|l|l|}
\hline
\textbf{Caratteristica} & \textbf{Tipologia} \\ \hline
Livello di alfabetizzazione & Alto \\ \hline
Titolo di studio & Laurea o Diploma superiore \\ \hline
Linguaggio nativo & Italiano \\ \hline
Abilità dattilografica & Alta \\ \hline
Alfabetizzazione informatica & Consolidata \\ \hline
\end{tabular}
\caption{Caratteristiche di conoscenza del Configuratore.}
\label{tab:conf_conoscenza}
\end{table}

\subsection{Caratteristiche di esperienza}
Il configuratore utilizza il sistema quotidianamente. Deve avere familiarità con altri sistemi di gestione per comprendere logiche di pianificazione e flussi di dati.

\begin{table}[H]
\centering
\begin{tabular}{|l|l|}
\hline
\textbf{Caratteristica} & \textbf{Tipologia} \\ \hline
Esperienza del compito & Alta \\ \hline
Esperienza sull'applicazione & Alta (utente abituale) \\ \hline
Esperienza dei sistemi informatici & Medio-Alta \\ \hline
Uso di altri sistemi & Frequente (ERP, gestionali) \\ \hline
\end{tabular}
\caption{Caratteristiche di esperienza del Configuratore.}
\label{tab:conf_esperienza}
\end{table}

\subsection{Caratteristiche fisiche e ambientali}
L'uso del sistema è obbligatorio e il compito è estremamente strutturato. L'interfaccia deve garantire l'accessibilità per sessioni di utilizzo prolungate.

\begin{table}[H]
\centering
\begin{tabular}{|l|l|}
\hline
\textbf{Caratteristica} & \textbf{Tipologia} \\ \hline
Età & 25 - 67 anni \\ \hline
Uso del sistema & Obbligatorio \\ \hline
Frequenza d'uso & Alta \\ \hline
Importanza del compito & Alta \\ \hline
Struttura del compito & Molto Alta \\ \hline
\end{tabular}
\caption{Caratteristiche fisiche e ambientali del Configuratore.}
\label{tab:conf_ambientali}
\end{table}

\subsection{Persona Narrativa}
\begin{itemize}
    \item \textbf{Background e Profilo Narrativo:} Dott. Marco Rossi, 45 anni. Coordinatore Logistico per l'Associazione Culturale responsabile della piattaforma. Vanta una formazione accademica umanistica affiancata ad una solida alfabetizzazione informatica gestionale, padroneggiando quotidianamente l'impiego di ERP, suite d'ufficio e fogli di calcolo avanzati.
    \item \textbf{Contesto Ambientale e d'Uso:} Opera prevalentemente dal proprio ufficio, su una workstation desktop con doppio monitor. Si trova cronicamente immerso in un ambiente frenetico, costantemente interrotto da chiamate telefoniche e urgenze organizzative, condizioni che rendono indispensabile un'interfaccia tollerante agli errori e ad alto recupero cognitivo.
    \item \textbf{Motivazioni e Obiettivi (Goals):} L'obiettivo psicologico dominante è il raggiungimento della tranquillità gestionale: desidera avere pieno controllo sulle tempistiche e sull'assegnazione dei volontari ai vari tour, azzerando le sovrapposizioni e massimizzando l'efficienza operativa per il successo dell'associazione.
    \item \textbf{Frustrazioni e Pain Points:} Insofferenza verso interfacce lente o proceduralmente farraginose. È frustrato da operazioni ripetitive (come l'inserimento manuale di eventi seriali) e dalla mancanza di micro-feedback visivi istantanei che confermino la robustezza o il salvataggio dei dati critici immessi.
    \item \textbf{Quote/Citazione rappresentativa:} \textit{``La tecnologia deve essere il mio braccio destro, non un ostacolo burocratico. Quando pianifico le attività, voglio interazioni nette e rassicuranti, che mi confermino immediatamente di avere la situazione in pugno.''}
\end{itemize}

\subsection{Obiettivi, azioni principali e differenze di interazione}
\textbf{Obiettivi principali:}
\begin{itemize}
    \item Pianificare e assegnare le visite garantendo copertura completa delle risorse umane disponibili.
    \item Gestire l'anagrafica di luoghi e tipologie di visita mantenendo il catalogo aggiornato.
    \item Monitorare lo stato avanzamento di ogni prenotazione, intervenendo prontamente su anomalie o conflitti di pianificazione.
\end{itemize}

\textbf{Azioni principali nel sistema:}
\begin{enumerate}
    \item Creazione e modifica di \textit{VisitType} (tipologie di tour) e \textit{Place} (luoghi).
    \item Revisione delle richieste pendenti nel pannello \textit{Visit Planning} e assegnazione di un volontario compatibile.
    \item Gestione del ciclo di vita delle visite: transizione degli stati da \textit{proposed} → \textit{confirmed} → \textit{effected} → \textit{complete}/\textit{cancelled}.
    \item Consultazione della dashboard amministrativa per statistiche aggregate su prenotazioni e volontari.
\end{enumerate}

\textbf{Differenze di interazione rispetto agli altri ruoli:}
Il Configuratore è l'unico utente con visibilità \emph{globale} su tutte le entità del sistema. A differenza del Volontario — che vede solo le proprie disponibilità e i tour assegnati — e del Fruitore — che opera esclusivamente sul proprio fascicolo prenotazioni — il Configuratore agisce su griglie dati dense (pattern CRUD) e flussi decisionali complessi. L'interfaccia gli richiede un \textit{information processing} di tipo \textbf{analitico-sequenziale}: deve leggere, confrontare e modificare molteplici record in un'unica sessione. Le conseguenti esigenze progettuali si traducono in: tabelle ordinabili e filtrabili, feedback sincrono sulle operazioni critiche, e controlli di validazione che prevengano conflitti di assegnazione.

\section{Volontario (Guida e fornitore di disponibilità)}
Il volontario mette a disposizione il proprio tempo per guidare i fruitori. Interagisce con il sistema principalmente per comunicare disponibilità e consultare i tour assegnati.

\subsection{Caratteristiche psicologiche e di conoscenza}
Lo stile cognitivo è orientato alla comunicazione. L'attitudine è positiva, ma condizionata dalla semplicità d'uso: la tecnologia non deve essere un ostacolo all'attività di volontariato.

\begin{table}[H]
\centering
\begin{tabular}{|l|l|}
\hline
\textbf{Caratteristica} & \textbf{Tipologia} \\ \hline
Stile cognitivo & Estroverso, Verbale, Intuitivo \\ \hline
Motivazione & Medio-Alta \\ \hline
Alfabetizzazione informatica & Media \\ \hline
Abilità dattilografica & Media \\ \hline
\end{tabular}
\caption{Caratteristiche psicologiche e di conoscenza del Volontario.}
\label{tab:vol_psico_conoscenze}
\end{table}

\subsection{Caratteristiche ambientali}
L'attività è facoltativa. Il contesto di utilizzo può essere dinamico (all'aperto), richiedendo un design dell'interfaccia chiaro e ad alto contrasto.

\begin{table}[H]
\centering
\begin{tabular}{|l|l|}
\hline
\textbf{Caratteristica} & \textbf{Tipologia} \\ \hline
Uso del sistema & Facoltativo \\ \hline
Frequenza d'uso & Medio-Bassa \\ \hline
Addestramento base & Nessuno (deve essere autoesplicativo) \\ \hline
Importanza del compito & Media \\ \hline
Uso di altri sistemi & Frequente (Social, Calendar) \\ \hline
\end{tabular}
\caption{Caratteristiche ambientali del Volontario.}
\label{tab:vol_ambientali}
\end{table}

\subsection{Persona Narrativa}
\begin{itemize}
    \item \textbf{Background e Profilo Narrativo:} Giulia Bianchi, 23 anni. Studentessa magistrale in Storia dell'Arte e Guida Volontaria appassionata. Possiede un'alfabetizzazione informatica orientata al \textit{mobile-first}: è un'utente espertissima di social network, app di messaggistica e calendari smart, ma risulta meno incline all'uso di elaborati cruscotti gestionali su PC.
    \item \textbf{Contesto Ambientale e d'Uso:} Interagisce con il sistema prevalentemente in mobilità assoluta (smartphone). Lo utilizza spessissimo all'aperto, in condizioni di luminosità avversa (es. forte luce solare) nei minuti immediatamente precedenti il ritrovo di un tour, o nei ritagli di tempo fra una lezione universitaria e l'altra.
    \item \textbf{Motivazioni e Obiettivi (Goals):} Cerca autorealizzazione personale condividendo la propria passione per l'arte senza l'ansia del lavoro burocratico. Il suo goal umano è riuscire a donare il proprio tempo libero alla comunità bilanciandolo armoniosamente con gli stringenti impegni accademici.
    \item \textbf{Frustrazioni e Pain Points:} Si avvilisce dinanzi a piattaforme non ottimizzate per \textit{touch-screen} o layout rigidi. È ostacolata dall'impossibilità di consultare agilmente i contatti del gruppo in caso di imprevisti last-minute o da flussi macchinosi per il ritiro e la gestione di una disponibilità data in precedenza.
    \item \textbf{Quote/Citazione rappresentativa:} \textit{``Il volontariato deve essere gratificante, non stressante. Ho bisogno di un sistema immediato, sul mio telefono, che mi dica chi sto per incontrare e quando, in modo chiaro e con pochissimi tap.''}
\end{itemize}

\subsection{Obiettivi, azioni principali e differenze di interazione}
\textbf{Obiettivi principali:}
\begin{itemize}
    \item Comunicare con precisione le proprie finestre di disponibilità per consentire all'amministrazione un'assegnazione ottimale.
    \item Consultare il calendario dei tour assegnati senza ambiguità su data, luogo e dettagli del gruppo.
    \item Minimizzare il tempo speso sul sistema: l'interazione deve essere rapida e concludibile in pochi passaggi.
\end{itemize}

\textbf{Azioni principali nel sistema:}
\begin{enumerate}
    \item Inserimento o aggiornamento delle disponibilità tramite il calendario interattivo (\textit{VolunteerAvailability}).
    \item Consultazione dell'elenco delle visite assegnate (\textit{volunteer.visits.past}) per ricevere dettagli logistici.
    \item Eventuale modifica o ritiro di una disponibilità precedentemente dichiarata.
\end{enumerate}

\textbf{Differenze di interazione rispetto agli altri ruoli:}
Il Volontario opera su un sottoinsieme circoscritto del sistema, privo di funzioni di back-office. Il suo modello di interazione privilegia la \textbf{manipolazione diretta} su oggetti visuali familiari (calendario) piuttosto che form testuali astratti. A differenza del Configuratore — che lavora prevalentemente su desktop in sessioni prolungate — il Volontario accede al sistema in micro-sessioni, spesso da smartphone, per aggiornamenti rapidi. Questo impone requisiti stringenti di \textit{mobile responsiveness}, riduzione del numero di tap per completare un'azione e feedback immediato che elimini ogni dubbio sul buon esito dell'operazione.

\section{Fruitore (Utente finale)}
Il fruitore utilizza la piattaforma per scoprire e prenotare le visite guidate. È l'utente che trae il beneficio culturale dal servizio.

\subsection{Caratteristiche psicologiche}
La motivazione è legata allo svago. Se il sito risulta complesso (alto carico cognitivo), l'utente potrebbe abbandonarlo facilmente in favore di altre fonti informative.

\begin{table}[H]
\centering
\begin{tabular}{|l|l|}
\hline
\textbf{Caratteristica} & \textbf{Tipologia} \\ \hline
Stile cognitivo & Concreto, Visivo, Intuitivo \\ \hline
Attitudine & Positiva (Ricerca di svago) \\ \hline
Motivazione & Moderata \\ \hline
\end{tabular}
\caption{Caratteristiche psicologiche del Fruitore.}
\label{tab:fruit_psicologiche}
\end{table}

\subsection{Caratteristiche di conoscenza e fisiche}
Il gruppo è estremamente eterogeneo (dai 18 agli 80 anni). Questo impone una particolare attenzione all'accessibilità (testi leggibili e compatibilità con screen reader).

\begin{table}[H]
\centering
\begin{tabular}{|l|l|}
\hline
\textbf{Caratteristica} & \textbf{Tipologia} \\ \hline
Livello di alfabetizzazione & Media (Popolazione generale) \\ \hline
Età & 18 - 80 anni \\ \hline
Alfabetizzazione informatica & Moderata-Bassa \\ \hline
Abilità dattilografica & Bassa / Non richiesta \\ \hline
\end{tabular}
\caption{Caratteristiche di conoscenza e fisiche del Fruitore.}
\label{tab:fruit_conoscenze_fisiche}
\end{table}

\subsection{Caratteristiche ambientali}
L'uso avviene in contesti di relax o in mobilità. L'importanza del compito è bassa (edonistica) e la navigazione è libera.

\begin{table}[H]
\centering
\begin{tabular}{|l|l|}
\hline
\textbf{Caratteristica} & \textbf{Tipologia} \\ \hline
Uso del sistema & Facoltativo \\ \hline
Frequenza d'uso & Occasionale \\ \hline
Importanza del compito & Bassa \\ \hline
Struttura del compito & Bassa \\ \hline
\end{tabular}
\caption{Caratteristiche ambientali del Fruitore.}
\label{tab:fruit_ambientali}
\end{table}

\subsection{Persona Narrativa}
\begin{itemize}
    \item \textbf{Background e Profilo Narrativo:} Antonio Ferrari, 62 anni. Insegnante di lettere in pensione, appassionato di storia ed escursionismo culturale. Presenta un'alfabetizzazione informatica basilare-moderata: sa navigare sul web, consultare la posta e acquistare biglietti per mezzi di trasporto, ma conserva una certa diffidenza generazionale verso meccanismi digitali iper-complessi e moduli intrusivi.
    \item \textbf{Contesto Ambientale e d'Uso:} Si approccia all'applicativo generalmente dal salotto di casa, nei momenti serali di relax, utilizzando un tablet o un laptop. L'ambiente è tranquillo, ma eventuali cali di vista dovuti all'età rendono le scelte tipografiche non accessibili particolarmente faticose sulla vista.
    \item \textbf{Motivazioni e Obiettivi (Goals):} Mira ad arricchire il proprio tempo libero e a organizzare uscite domenicali serene in compagnia di familiari o amici (spesso nipoti). L'obiettivo psicologico è acquisire la sicurezza totale di aver bloccato il proprio posto, rassicurando sé e i propri accompagnatori senza brutte sorprese all'arrivo.
    \item \textbf{Frustrazioni e Pain Points:} Prova forte ansia da \textit{checkout} informatico e diffidenza verso la condivisione esagerata di dati anagrafici prima di poter prendere visione delle tariffe. È infastidito da font minuti e contrasti cromatici labili (es. testo grigio tenue su sfondo bianco). L'incertezza sullo stato finale (es. ``la prenotazione è andata a buon fine o no?'') rappresenta il suo principale deterrente.
    \item \textbf{Quote/Citazione rappresentativa:} \textit{``Esplorare la mia città in compagnia deve essere un piacere fin dal primo click. Non voglio dover essere un ingegnere per assicurarmi quattro posti liberi e per accertarmi che la mia prenotazione sia realmente andata a buon fine.''}
\end{itemize}

\subsection{Obiettivi, azioni principali e differenze di interazione}
\textbf{Obiettivi principali:}
\begin{itemize}
    \item Individuare rapidamente un tour di interesse consultando il catalogo con filtri pertinenti (data, luogo, lingua).
    \item Completare la prenotazione con il minimo sforzo cognitivo, ricevendo conferma esplicita e inequivocabile dell'avvenuta iscrizione.
    \item Poter gestire autonomamente la propria prenotazione (consultazione, disdetta) senza necessità di supporto esterno.
\end{itemize}

\textbf{Azioni principali nel sistema:}
\begin{enumerate}
    \item Navigazione della \textit{Landing Page} e ricerca/filtro delle visite disponibili.
    \item Compilazione del modulo di richiesta visita (\textit{Registration}) specificando numero di partecipanti, lingua e note aggiuntive.
    \item Download del biglietto PDF con QR code a conferma dell'iscrizione.
    \item Cancellazione di una prenotazione esistente dalla propria dashboard (\textit{user.visits.past}).
\end{enumerate}

\textbf{Differenze di interazione rispetto agli altri ruoli:}
Il Fruitore rappresenta l'utente con il \textbf{minor controllo sul sistema} e, al tempo stesso, il \textbf{più alto rischio di abbandono}: a differenza del Configuratore e del Volontario — che utilizzano la piattaforma per obbligo o scelta consapevole — il Fruitore può rinunciare in qualsiasi momento a favore di alternative più semplici. Il suo modello cognitivo è \textbf{visivo-intuitivo}: si aspetta che le funzioni principali siano immediatamente evidenti (\textit{affordance}) senza richiedere alcun addestramento preliminare. L'ampiezza demografica del gruppo (18–80 anni) con gradi eterogenei di alfabetizzazione informatica impone il rispetto rigoroso dei principi di accessibilità: testo leggibile, contrasto elevato, aree di tocco generose e feedback di stato sempre espliciti. Queste differenze si riflettono direttamente nelle scelte di redesign discusse nei capitoli successivi.